% Propietario conservado: /Users/ruben/Documents/New project/output/REVISION_ARGUMENTAL_APERTURAS_CIERRES_20260915/III/spanish_source/manuscrito/sections/11_nucleo_finito.tex
\section{Énumération exhaustive et détermination finie des opérateurs de transport}
\label{iii:iii:nucleo-finito}

Cet appendice développe deux opérations finies avec des données de départ
différentes. La première énumère les conditions initiales des deux curseurs
et construit leurs émissions au moyen de la récurrence arithmétique déclarée.
La seconde détermine les opérateurs de transport à partir du registre
régional des transitions. La séparation de leurs domaines permet de préciser
ce que produit chaque démonstration : l’énumération construit le catalogue des
émissions ; la reconstruction matricielle détermine les opérateurs compatibles
avec le registre. Le théorème de reconstruction n’est pas employé comme démonstration
de la production antérieure de ce registre.

Les primitives de la première opération sont le support $I_9^2$, les quatre
orientations cardinales, le sélecteur tritique, la lecture antérieure au
déplacement, les inversions d’orientation et la règle d’émission.
Leurs définitions sont réunies ci-dessous. Celles de la seconde sont les
quatre matrices de transitions de~\eqref{iii:iii:nf-registro}.
Leur reconstruction unique est démontrée dans le
théorème~\ref{iii:iii:nf-unicidad} ; le calendrier qui en résulte est exprimé
dans~\eqref{iii:iii:nf-calendario}. Ces déterminations conservent
les données de départ déclarées pour chaque opération.

\subsection{Domaine complet des conditions initiales}

Soit $I_9=\{0,\ldots,8\}$ et soit
\[
 \mathcal D=\{N,E,S,O\},\qquad
 v_N=(-1,0),\quad v_E=(0,1),\quad
 v_S=(1,0),\quad v_O=(0,-1).
\]
Une condition initiale de curseur est $s=(i,j,d)\in
\mathcal S=I_9^2\times\mathcal D$. Les deux curseurs sont indépendants
dans le domaine initial
\begin{equation}
 \Omega=\mathcal S_+\times\mathcal S_\times,
 \qquad |\mathcal S|=9^2\cdot4=324,
 \qquad |\Omega|=324^2=104\,976.
 \label{iii:iii:nf-dominio}
\end{equation}
L’indice $i$ correspond à l’étiquette positive $i+1$ d’APP. Pour
des entiers positifs, la réduction positive est $\rho_9(n)=1+[(n-1)]_9$,
où $[a]_b$ désigne le représentant dans $\{0,\ldots,b-1\}$.
Les évaluations résiduelles des deux feuillets sont
\[
 \sigma(i,j)=\rho_9(i+j+2),\qquad
 \varpi(i,j)=\rho_9((i+1)(j+1)).
\]
Le lecteur multiplicatif fini utilise
\begin{equation}
 \begin{array}{c|rrrrrrrrr}
 a&1&2&3&4&5&6&7&8&9\\ \hline
 \ell_9(a)&0&1&0&2&5&0&4&3&0
 \end{array}.
 \label{iii:iii:nf-log}
\end{equation}
Sur les unités, $\ell_9$ est l’exposant discret en base $2$.
Sur $3,6,9$, la valeur $0$ constitue l’extension déclarée de
l’observable d’émission ; l’appartenance à ce secteur demeure dans le
registre de trajectoire.

Pour $1\le t\le54$, on définit
\[
 r_t=1+[(t-1)]_9,
 \qquad
 \tau_t=
 \begin{cases}
 +1,&r_t\in\{1,4,7\},\\
 -1,&r_t\in\{2,5,8\},\\
 0,&r_t\in\{3,6,9\}.
 \end{cases}
\]
Chaque curseur conserve sa position $z_t$ immédiatement avant le pas
$t$ et son vecteur d’orientation $v_t$. Le curseur additif est actif
lorsque $\tau_t=+1$ ; le multiplicatif, lorsque $\tau_t=-1$. Pour chacun
d’entre eux, avec l’indicateur d’activité $a_t$, la mise à jour est
\begin{equation}
 z_{t+1}=[z_t+a_tv_t]_9,
 \qquad
 v_{t+1}=
 \begin{cases}
 -v_t,&t\in\{27,54\},\\
 v_t,&t\notin\{27,54\}.
 \end{cases}
 \label{iii:iii:nf-actualizacion}
\end{equation}
La lecture du pas $t$ est effectuée en $z_t$, avant le déplacement.
L’événement neutre conserve les deux positions. Cette convention fixe
sans ambiguïté les extrémités des fenêtres et l’effet des
inversions cardinales.

Dans les fenêtres $W_m=\{9m-8,\ldots,9m\}$, $1\le m\le6$, soient
\begin{align}
 S_m(s_+)&=\sum_{\substack{t\in W_m\\\tau_t=+1}}
                 \sigma(z_t^+),&
 E_m(s_\times)&=\sum_{\substack{t\in W_m\\\tau_t=-1}}
                 \ell_9(\varpi(z_t^\times)),
 \label{iii:iii:nf-firmas}\\
 s_m&=[S_m]_{10},&e_m&=[E_m]_{10}.\notag
\end{align}
Dans chaque fenêtre, il y a trois événements de chaque type. Les signatures $s$ et $e$
sont les six composantes de ces résidus, non les entiers $S_m,E_m$
antérieurs à leur réduction.

\subsection{Recensement des signatures et des émissions}

L’émission d’une paire de conditions initiales est déterminée par
\begin{equation}
 U_m=100s_m+10[s_m+e_m]_{10}+[3s_m+5e_m+1]_{10},
 \qquad 1\le m\le6.
 \label{iii:iii:nf-emision}
\end{equation}
Le terme final $1$ est $[7\cdot3]_{10}$, correspondant aux trois
événements neutres. Les coefficients de cette règle font partie de la
récurrence déclarée. Les constantes réelles étudiées ultérieurement
n’interviennent ni dans l’énumération de son domaine ni dans son évaluation.

\begin{proposition}[Factorisation injective du catalogue fini]
\label{iii:iii:nf-inyectividad}
L’émission distingue chaque paire de signatures $(s,e)$. Il y a $18$ signatures
additives et $26$ multiplicatives ; leurs paires produisent $468$ émissions.
Les multiplicités de ces émissions sont $144$, $432$ et $1944$,
avec $432$, $18$ et $18$ émissions respectivement.
\end{proposition}

\begin{proof}
Le chiffre des centaines de $U_m$ restitue $s_m$. Si $b_m$ est son chiffre des
dizaines, alors $e_m=[b_m-s_m]_{10}$. Le chiffre des unités vérifie
la troisième congruence de \eqref{iii:iii:nf-emision}. Par conséquent, le
couplage est injectif.

Le recensement s’obtient en parcourant, dans l’ordre cartésien complet,
$i=0,\ldots,8$, $j=0,\ldots,8$ et $d=N,E,S,O$, et en appliquant les $54$
pas de \eqref{iii:iii:nf-actualizacion}. Les $324$ conditions additives
se répartissent entre les $18$ signatures suivantes, chacune avec $18$
préimages. Dans les tableaux, une suite de six chiffres représente
les six composantes ordonnées d’une signature.
\begin{center}
\begin{tabular}{ccc}
$122564$&$122898$&$212645$\\
$212988$&$221456$&$221889$\\
$456221$&$465898$&$546988$\\
$564122$&$645212$&$654889$\\
$889221$&$889654$&$898122$\\
$898465$&$988212$&$988546$
\end{tabular}
\end{center}
Pour le curseur multiplicatif, les signatures et leurs multiplicités sont
\begin{center}
\begin{tabular}{rr@{\qquad}rr}
Signature&Multiplicité&Signature&Multiplicité\\ \hline
$000000$&$108$&$555555$&$24$\\
$177393$&$8$&$555717$&$8$\\
$177555$&$8$&$555771$&$8$\\
$177771$&$8$&$555933$&$8$\\
$339555$&$8$&$717555$&$8$\\
$339771$&$8$&$717717$&$8$\\
$339933$&$8$&$717933$&$8$\\
$393177$&$8$&$771177$&$8$\\
$393393$&$8$&$771339$&$8$\\
$393555$&$8$&$771555$&$8$\\
$555177$&$8$&$933339$&$8$\\
$555339$&$8$&$933555$&$8$\\
$555393$&$8$&$933717$&$8$
\end{tabular}
\end{center}
Ces tableaux résultent des sommes finies \eqref{iii:iii:nf-firmas}, dont le
domaine, l’ordre de lecture et les inversions sont déjà spécifiés. Leurs sommes
de multiplicités sont $18\cdot18=324$ et
$24\cdot8+24+108=324$. Le domaine cartésien
\eqref{iii:iii:nf-dominio} permet de réaliser n’importe quelle paire de signatures.
L’injectivité fournit $18\cdot26=468$ émissions et les
multiplicités
\[
 (18\cdot24,\ 18\cdot8)=(432,144),\qquad
 (18,\ 18\cdot24)=(18,432),\qquad
 (18,\ 18\cdot108)=(18,1944),
\]
où chaque paire indique le nombre d’émissions et la multiplicité de sa fibre.
Enfin,
\[
 432\cdot144+18\cdot432+18\cdot1944=104\,976.
\]
L’énumération épuise ainsi le domaine complet des conditions initiales.
\end{proof}

La réduction ternaire orientée est
\[
 \Psi(U)=([-U_1]_3,\ldots,[-U_6]_3).
\]
Appliquée aux $18\cdot26$ émissions des tableaux précédents,
elle produit le recensement des fibres suivant. Ici, $r$ compte les émissions
distinctes, non les conditions initiales.
\begin{equation}
 \begin{array}{c|rrrrrr}
 r&1&2&3&4&5&6\\ \hline
 \#\{w:\#\Psi^{-1}(w)=r\}&132&66&18&3&6&18
 \end{array}.
 \label{iii:iii:nf-censo-ternario}
\end{equation}
La somme de la deuxième ligne est $243$, tandis que sa somme pondérée
par $r$ est $468$. On distingue ainsi l’image de $243$ mots
de l’espace ambiant $\mathbb F_3^6$, qui contient $729$ mots. Le certificat
joint conserve toutes les signatures avec leurs préimages, les $468$
émissions et leur réduction ternaire ; les formules précédentes fournissent
une procédure complète pour les reconstruire.

\subsection{Registre des transitions et reconstruction des relèvements}

Dans cette partie, on fixe le registre régional du propriétaire cité.
Soient $b_0^\chi,\ldots,b_4^\chi\in\mathbb F_3^6$, avec
$\chi\in\{\mathsf C,\mathsf P,\mathsf A\}$, ses trois suites.
Le registre organise les transitions $b_0\to b_1\to b_2$ dans
$X_0,Y_0$ et les transitions $b_2\to b_3\to b_4$ dans $X_1,Y_1$,
avec des vecteurs lignes et le même ordre régional :
\begingroup\small
\begin{equation}
\begin{aligned}
X_0&=\begin{pmatrix}
0&1&0&2&1&1\\0&1&2&2&2&2\\2&0&1&1&0&1\\
1&2&1&2&2&1\\1&2&1&2&0&0\\1&1&2&2&0&2
\end{pmatrix},&
Y_0&=\begin{pmatrix}
0&1&2&2&2&2\\0&1&0&2&1&1\\1&2&1&2&2&1\\
1&0&2&0&1&1\\1&1&2&2&0&2\\1&2&1&0&2&0
\end{pmatrix},\\[1ex]
X_1&=\begin{pmatrix}
0&1&0&2&1&1\\0&0&2&1&1&1\\1&0&2&0&1&1\\
0&1&2&2&2&2\\1&2&1&0&2&0\\0&1&0&2&1&0
\end{pmatrix},&
Y_1&=\begin{pmatrix}
0&0&2&1&1&1\\1&1&0&2&2&1\\0&1&2&2&2&2\\
1&0&2&0&1&1\\0&1&0&2&1&0\\0&1&0&2&0&0
\end{pmatrix}.
\end{aligned}
\label{iii:iii:nf-registro}
\end{equation}
\endgroup
La production précédente s’exprime chez son propriétaire au moyen de la
composition de sélection, de transport et de mise à jour :
\[
 x_{j+1}^{\chi}
 =\mathcal U_{9j+9}\circ\cdots\circ\mathcal U_{9j+1}(x_j^{\chi}),
 \qquad b_j^\chi=\Pi_6(x_j^\chi),
 \qquad\mathcal U_t=\mathsf{Upd}_t\circ\mathsf{Tra}_t\circ\mathsf{Sel}_t.
\]
Le résultat suivant part du registre \eqref{iii:iii:nf-registro}.
Sa démonstration permet de reconstruire les opérateurs et de vérifier la
compatibilité des transitions ; elle n’attribue pas à l’inversion matricielle
la production coordonnée des états $x_j^\chi$.

\begin{proposition}[Unicité sur le registre régional]
\label{iii:iii:nf-unicidad}
Les systèmes $X_aL_a=Y_a$, $a=0,1$, ont une solution unique sur
$\mathbb F_3$. Les deux solutions sont
\begingroup\small
\begin{equation}
L_0=\begin{pmatrix}
2&2&2&1&2&1\\2&1&2&2&1&1\\1&0&0&1&0&2\\
0&0&1&1&1&0\\2&2&2&0&2&1\\2&1&2&1&0&0
\end{pmatrix},\qquad
L_1=\begin{pmatrix}
2&1&2&0&2&2\\1&2&1&1&2&0\\1&2&1&1&1&2\\
0&1&0&0&2&1\\2&0&2&1&0&1\\0&2&2&2&1&1
\end{pmatrix}.
\label{iii:iii:nf-lifts}
\end{equation}
\endgroup
\end{proposition}

\begin{proof}
L’élimination dans $\mathbb F_3$ donne $\det X_0=1$ et
$\det X_1=2=-1$. Chaque application $L\mapsto X_aL$ est donc
un automorphisme de l’espace des matrices de dimension $36$ ; dans une
vectorisation par colonnes, elle est représentée par $I_6\otimes X_a$.
On obtient $L_a=X_a^{-1}Y_a$. La multiplication des matrices de
\eqref{iii:iii:nf-lifts} par les matrices $X_a$ respectives reproduit les
deux matrices $Y_a$ de \eqref{iii:iii:nf-registro} ; existence et unicité
sont établies sur ce registre.
\end{proof}

Les trois suites résultantes, avec une séparation explicite en blocs,
sont
\begin{align*}
 B^{\mathsf C}&=010211\mid012222\mid010211\mid002111\mid110221,\\
 B^{\mathsf P}&=201101\mid121221\mid102011\mid012222\mid102011,\\
 B^{\mathsf A}&=121200\mid112202\mid121020\mid010210\mid010200.
\end{align*}
Chaque bloc conserve ses six positions. La concaténation n’est pas
identifiée à un entier en base dix.

\subsection{Exhaustivité du calendrier et stabilité de la reconstruction}

Un calendrier binaire est $c=(c_0,c_1,c_2,c_3)\in\{0,1\}^4$.
On définit le nombre de transitions incompatibles avec le registre par
\begin{equation}
 d(c)=\sum_{\chi\in\{\mathsf C,\mathsf P,\mathsf A\}}
       \sum_{j=0}^{3}
       \mathbf1\{b_j^\chi L_{c_j}\ne b_{j+1}^\chi\}.
 \label{iii:iii:nf-calendario}
\end{equation}
Les multiplications des suites précédentes par
\eqref{iii:iii:nf-lifts} donnent
\[
 d(c)=3\,d_{\rm H}(c,(0,0,1,1)),
\]
où $d_{\rm H}$ est la distance de Hamming : à chacune des
quatre positions, l’opérateur alternatif ne satisfait pas les trois
transitions régionales. Le recensement complet est
\begin{center}
\begin{tabular}{cr@{\qquad}cr@{\qquad}cr@{\qquad}cr}
$c$&$d(c)$&$c$&$d(c)$&$c$&$d(c)$&$c$&$d(c)$\\ \hline
$0000$&$6$&$0100$&$9$&$1000$&$9$&$1100$&$12$\\
$0001$&$3$&$0101$&$6$&$1001$&$6$&$1101$&$9$\\
$0010$&$3$&$0110$&$6$&$1010$&$6$&$1110$&$9$\\
$0011$&$0$&$0111$&$3$&$1011$&$3$&$1111$&$6$
\end{tabular}
\end{center}
Par conséquent, $0011$ est l’unique calendrier qui reproduit les
trois suites complètes.

La sensibilité aux entrées des matrices admet également une
preuve exhaustive sans approximations. Une modification unitaire
remplace $L_a$ par $L_a+\lambda E_{ij}$, avec
$a\in\{0,1\}$, $1\le i,j\le6$ et
$\lambda\in\mathbb F_3\setminus\{0\}$. Il y a
$2\cdot6\cdot6\cdot2=144$ modifications de ce type. Comme
$X_a$ est inversible,
\[
 X_a(L_a+\lambda E_{ij})-Y_a
 =\lambda X_aE_{ij}\ne0.
\]
Toute modification détruit au moins une des transitions du
registre correspondant. Le certificat exécutable énumère en outre
les $144$ modifications et conserve, pour chacune, le nombre de
transitions non satisfaites.

La propriété démontrée est la rigidité des opérateurs par rapport au
registre régional fixé. Pour reproduire une construction antérieure
du registre lui-même, l’action coordonnée de ses opérateurs
de sélection, de transport et de mise à jour est nécessaire. Cette dépendance reste
séparée du théorème de rigidité, du catalogue fini des émissions et
des réalisations analytiques ultérieures.

\subsection{Certificat de reproduction et résidence des données}

Le programme \texttt{verificar\_nucleo\_finito.py} parcourt les $324$
conditions de chaque curseur et forme le produit des classes résultantes.
Le catalogue documentaire est chargé après cette construction pour
comparer les émissions et leurs multiplicités. Dans la seconde partie,
les quatre matrices de transitions sont extraites de la copie littérale
du propriétaire ; les solutions sont calculées par élimination exacte
modulo $3$ avant de les comparer aux matrices imprimées. Le reçu
\texttt{NUCLEO\_FINITO.json} conserve les domaines, les préimages,
les émissions, les $16$ calendriers et les $144$ modifications.

La vérification utilise des entiers et des résidus exacts. Elle ne reçoit ni valeurs
métrologiques ni développements réels des constantes. Sa portée
est celle des opérations finies et des données de départ spécifiées
dans cet appendice. Les démonstrations de compatibilité à profondeur
illimitée et les identifications physiques conservent leurs domaines
et leurs obligations d’exposition respectifs.
